\medskip

\emph{Cet exercice est un questionnaire à choix multiples. Pour chacune des questions suivantes, une seule des quatre réponses proposées est exacte.}

\medskip

\emph{Une réponse incorrecte, une réponse multiple ou l'absence de réponse à une question ne rapporte ni n'enlève de point. Pour répondre, indiquer sur la copie le numéro de la question et la lettre de la réponse choisie.\\
Aucune justification n'est demandée.}

\medskip

\begin{enumerate}
\item Soit $f$ la fonction définie sur $\R$ par

\[f(x) =\dfrac{x}{\text{e}^x}\]

On suppose que $f$ est dérivable sur $\R$ et on note $f'$ sa fonction dérivée. 

\begin{center}
\begin{tabularx}{\linewidth}{*{2}{X}}
\textbf{a.~~}$f'(x) = \text{e}^{-x}$& \textbf{b.~~} $f'(x) = x\text{e}^{-x}$\\
\textbf{c.~~}$f'(x) = (1- x)\text{e}^{-x}$& \textbf{d.~~} $f'(x) = (1+x) \text{e}^{-x}$
\end{tabularx}
\end{center}

\item Soit $f$ une fonction deux fois dérivable sur l'intervalle $[-3~;~1]$. On donne ci-dessous la représentation graphique de sa fonction dérivée seconde $f''$.

\begin{center}
\psset{xunit=1.5cm,yunit=0.75cm}
\begin{pspicture*}(-3.5,-3.5)(1.5,3.5)
\psgrid[gridlabels=0pt,subgriddiv=1,gridwidth=0.2pt]
\psaxes[linewidth=1.25pt]{->}(0,0)(-3.5,-3.5)(1.5,3.5)
\psplot[plotpoints=2000,linewidth=1.25pt,linecolor=red]{-3}{1}{x 3 exp x dup mul 3 mul add x sub 3 sub}
\end{pspicture*}

\end{center}

On peut alors affirmer que:

\begin{center}
\begin{tabularx}{\linewidth}{*{2}{X}}
\textbf{a.~~}La fonction $f$ est convexe sur l'intervalle $[-1~;~1]$&\textbf{b.~~} La fonction $f$ est concave sur l'intervalle $[- 2~;~0]$\\
\textbf{c.~~}La fonction $f'$ est décroissante sur l'intervalle $[-2~;~0]$&\textbf{d.~~}
La fonction $f'$ admet un maximum en 

$x = -1$
\end{tabularx}
\end{center}


\item On considère la fonction $f$ définie sur $\R$ par:

\[f(x) = x^3\text{e}^{-x^2}\]

Si $F$ est une primitive de $f$ sur $\R$, 

\begin{center}
\renewcommand\arraystretch{1.9}
\begin{tabularx}{\linewidth}{*{2}{X}}
\textbf{a.~~}$F(x) = - \dfrac16\left(x^3 + 1\right)\text{e}^{-x^2}$ &
\textbf{b.~~}$F(x) = - \dfrac14 x^4\text{e}^{-x^2}$\\
\textbf{c.~~}$F(x)=-\dfrac12\left(x^2 + 1\right)\text{e}^{-x^2}$&
\textbf{d.~~}$F(x)= x^2\left(3 - 2x^2\right)\text{e}^{-x^2}$
\end{tabularx}
\renewcommand\arraystretch{1}
\end{center}

\item Que vaut :

\[\displaystyle\lim_{x \to + \infty} \dfrac{\text{e}^x + 1}{\text{e}^x - 1}\]

\begin{center}
\begin{tabularx}{\linewidth}{*{2}{X}}
\textbf{a.~~} $-1$&\textbf{b.~~}  1\\
\textbf{c.~~} $+ \infty$&\textbf{d.~~} n'existe pas
\end{tabularx}
\end{center}

\item On considère la fonction $f$ définie sur $\R$ par $f(x) = \text{e}^{2x +1}$.

La seule primitive $F$ sur $\R$ de la fonction $f$ telle que $F(0) = 1$ est la fonction :

\begin{center}
\begin{tabularx}{\linewidth}{*{2}{X}}
\textbf{a.~~}$x \longmapsto  2\text{e}^{2x+1} - 2\text{e} + 1$&\textbf{b.~~}$x \longmapsto  2\text{e}^{2x+1} - \text{e}$\\
\textbf{c.~~} $x \longmapsto \dfrac12\text{e}^{2x+1} - \dfrac12 \text{e} + 1$&\textbf{d.~~}$x \longmapsto \text{e}^{x^2 + x}$
\end{tabularx}
\end{center}

\item ~

\begin{minipage}{0.46\linewidth}
Dans un repère, on a tracé ci-contre la courbe représentative d'une fonction $f$ définie et deux fois dérivable sur $[- 2 ~;~4]$
\end{minipage} \hfill \begin{minipage}{0.48\linewidth}
\psset{unit=0.9cm}
\begin{pspicture}(-3,-2)(5,3)
\psgrid[gridlabels=0pt,subgriddiv=1,gridwidth=0.2pt]
\psaxes[linewidth=1.25pt,Dx=10,Dy=10]{->}(0,0)(-3,-2)(5,3)
\psaxes[linewidth=1.25pt](0,0)(1,0)(0,1)
\psplot[plotpoints=2000,linewidth=1.25pt,linecolor=red]{-2}{4}{x 3 exp 3 div x dup mul sub 8 3 div add  3 mul 8 div}
\uput[d](4.6,0){$x$}\uput[l](0,2.6){$y$}
\end{pspicture}

\end{minipage}

Parmi les courbes suivantes, laquelle représente la fonction $f''$,
dérivée seconde de $f$ ? 

\begin{center}
\begin{tabularx}{\linewidth}{*{2}{X}}
\textbf{a.~~} \psset{unit=0.6cm}
\begin{pspicture}(-3,-3)(5,3)
\psgrid[gridlabels=0pt,subgriddiv=1,gridwidth=0.2pt]
\psaxes[linewidth=1.25pt,Dx=10,Dy=10]{->}(0,0)(-3,-3)(5,3)
\psplot[plotpoints=2000,linewidth=1.25pt,linecolor=blue]{-2}{4}{x 1 sub }
\psaxes[linewidth=1.25pt](0,0)(1,0)(0,1)
\uput[d](4.6,0){$x$}\uput[l](0,2.6){$y$}
\end{pspicture}&\textbf{b.~~} \psset{unit=0.6cm}
\begin{pspicture}(-3,-3)(5,3)
\psgrid[gridlabels=0pt,subgriddiv=1,gridwidth=0.2pt]
\psaxes[linewidth=1.25pt,Dx=10,Dy=10]{->}(0,0)(-3,-3)(5,3)
\psplot[plotpoints=2000,linewidth=1.25pt,linecolor=blue]{-2}{4}{1 x sub}
\psaxes[linewidth=1.25pt](0,0)(1,0)(0,1)
\uput[d](4.6,0){$x$}\uput[l](0,2.6){$y$}
\end{pspicture}\\
\textbf{c.~~} \psset{unit=0.6cm}
\begin{pspicture}(-3,-2)(5,4)
\psgrid[gridlabels=0pt,subgriddiv=1,gridwidth=0.2pt]
\psaxes[linewidth=1.25pt,Dx=10,Dy=10]{->}(0,0)(-3,-2)(5,4)
\psplot[plotpoints=2000,linewidth=1.25pt,linecolor=blue]{-2}{4}{x 3  div x 2 sub mul}
\psaxes[linewidth=1.25pt](0,0)(1,0)(0,1)
\uput[d](4.6,0){$x$}\uput[l](0,3.6){$y$}
\end{pspicture}& \textbf{d.~~}
\psset{unit=0.6cm}
\begin{pspicture}(-3,-2)(5,4)
\psgrid[gridlabels=0pt,subgriddiv=1,gridwidth=0.2pt]
\psaxes[linewidth=1.25pt,Dx=10,Dy=10]{->}(0,0)(-3,-2)(5,4)
\psplot[plotpoints=2000,linewidth=1.25pt,linecolor=blue]{-1.8}{3.8}{x 1  sub dup mul  2 div}
\psaxes[linewidth=1.25pt](0,0)(1,0)(0,1)
\uput[d](4.6,0){$x$}\uput[l](0,3.6){$y$}
\end{pspicture}
\end{tabularx}
\end{center}
\end{enumerate}

\bigskip
